%%%%%
%@TheDoctorRAB
%%%%%
%
%%%%%
%standard white paper/preproposal format
%%%%%
%
%%%%%
%PDF tagging accessible file
%compile with pdflatex
%run $show-pdf-tags --xml filename.pdf
%%%%%
%
%%%%%
%REFERENCES
%neup.bst - numbered citations in order of appearance, short author list with et al in reference section
%nsf.bst - numbered citations in order of appearance, full author list in references section
%standard.bst - citations with author last name with et al for more than 2 authors; full author list in references section
%ans.bst is for ANS only. 
%
%author = {Lastname, Firstname and Lastname, Firstname and Lastname, Firstname} for all bst formats
%bst renders the author list itself
%
%author = {{Nuclear Regulatory Commission}} if the author is an organization, institution, etc., and not people
%
%title = {{}} for all
%
%for all - use \citep{-} - [1] or (Borrelli, 2021) in the text
%standard.bst \cite{-} - Borrelli (2021) in the text
%standard.bst lists references alphabetically
%the rest list numerically
%%%%%


%%%%% enable tagging for accessibility
%turn off if using titlesec,titletoc
\RequirePackage{pdfmanagement-testphase}
\DocumentMetadata{
    tagging=on,
    pdfversion=2.0,
    pdfstandard=ua-2,
    lang=en-US,
    debug={log=all},
    testphase={latest,tagpdf},
    tagging-setup={math/setup={mathml-AF,mathml-SE}}
}
\tagpdfsetup{math/alt/use}
%%%%%


%%%%% general 
\documentclass[11pt,a4paper]{article}
\usepackage[lmargin=1in,rmargin=1in,tmargin=1in,bmargin=1in]{geometry}
\usepackage[pagewise]{lineno} %linenumbering restarts each page
%\usepackage[running]{lineno} %linenumbering continuous
\usepackage{mmap} %helps make pdfs copy and paste
\usepackage{setspace}
\usepackage{ulem} %strikethrough - do not \sout{\cite{}}
\usepackage[pdftex,dvipsnames]{xcolor,colortbl} %change font color
\usepackage{graphicx}
\usepackage{tablefootnote}
\usepackage{footnotehyper}
\usepackage{float}
%\usepackage{mypythonhighlight,verbatim}
%\usepackage{subfig}
\usepackage[yyyymmdd]{datetime} %date format
\renewcommand{\dateseparator}{.}
%\graphicspath{{$TEXIMG/}} %path to graphics
\setcounter{secnumdepth}{5} %set subsection to nth level
\usepackage{needspace}
\usepackage[stable,hang,flushmargin]{footmisc} %footnotes in section titles and no indent; standard.bst
\usepackage{enumext}
%\usepackage[inline]{enumitem} %https://www.tug.org/texlive//Contents/live/texmf-dist/doc/latex-dev/latex-lab/latex-lab-enumitem.pdf
\usepackage{boldline}
\usepackage{makecell}
\usepackage{booktabs}
\usepackage{amssymb}
\usepackage{gensymb}
\usepackage{amsmath,nicefrac}
\usepackage{physics}
\usepackage{lscape}
\usepackage{array}
\usepackage{chngcntr}
\usepackage{sectsty}
\usepackage{textcomp}
\usepackage{lastpage}
\usepackage{xargs} %for \newcommandx
\usepackage[colorinlistoftodos,prependcaption,textsize=tiny]{todonotes} %makes colored boxes for commenting
\usepackage{soul}
\usepackage{color}
\usepackage{marginnote}
\usepackage[figure,table]{totalcount}
\usepackage{parskip}
\usepackage{multicol}
%%%%%


%%%%% links and tags
\usepackage{hyperref}
\hypersetup{
    colorlinks,
    linkcolor=black,
    citecolor=black,
    urlcolor=blue,
    pdftitle={Preproposal}, %turn off if tagging off
    pdfauthor={R. A. Borrelli}, %turn off if tagging off
    pdfkeywords={preproposal}, %turn off if tagging off
    pdflanguage={en-US}, %turn off if tagging off
    pdfsubject={preproposal}, %turn off if tagging off
    pdfdisplaydoctitle=true  %turn off if tagging off
} 
\usepackage[capitalise]{cleveref} %must be loaded after hyperref
%%%%%


%%%%% tikz
\usepackage{pgf}
\usepackage{tikz} % required for drawing custom shapes
\usetikzlibrary{shapes,arrows,automata,trees}
%%%%%


%%%%% fonts
\usepackage{times}
%arial - uncomment next two lines
%\usepackage{helvet}
%\renewcommand{\familydefault}{\sfdefault}
%%%%%


%%%%% references
%\usepackage[round,semicolon]{natbib} %for (Borrelli 2021; Clooney 2019) - standard.bst 
\usepackage[numbers,sort&compress]{natbib} %for [1-3] - nsf.bst, neup.bst
\setlength{\bibsep}{7pt} %sets space between references
%\renewcommand{\bibsection}{} %suppresses large 'references' heading
%\renewcommand\bibpreamble{\vspace{\baselineskip}} %sets spacing after heading if not using default references heading
%%%%%


%%%%% tables and figures
\usepackage{longtable} %need to put label at top under caption then \\ - use spacing
\usepackage{tablefootnote}
\usepackage{tabularx}
\usepackage{multirow}
\usepackage{tabto} %general tabbed spacing
\usepackage{pdfpages}
\usepackage{wrapfig} %wraps figures around text
\setlength{\intextsep}{0.00mm}
\setlength{\columnsep}{1.00mm}
\usepackage[singlelinecheck=false,labelfont=bf]{caption}
\usepackage{subcaption}
\captionsetup[table]{justification=centering,skip=5pt,labelformat={default},labelsep=period,name={Table}} %sets a space after table caption
\captionsetup[figure]{justification=justified,skip=5pt,labelformat={default},labelsep=period,name={Figure}} %sets space above caption, 'figure' format
\captionsetup[wrapfigure]{justification=centering,aboveskip=0pt,belowskip=0pt,labelformat={default},labelsep=period,name={Fig.}} %sets space above caption, 'figure' format
\captionsetup[wraptable]{justification=centering,aboveskip=0pt,belowskip=0pt,labelformat={default},labelsep=period,name={Table}} %sets space above caption, 'figure' format
%%%%%


%%%%% watermark
%\usepackage[firstpage,vpos=0.63\paperheight]{draftwatermark}
%\SetWatermarkText{\shortstack{DRAFT\\do not distribute}}
%\SetWatermarkScale{0.20}
%%%%%


%%%%% cross referencing files
%\usepackage{xr} %for revisions - will cross reference from one file to here
%\externaldocument{/path/to/auxfilename} %aux file needed
%%%%%


%%%%% toc and glossaries
\usepackage[toc,title]{appendix}
\usepackage[acronym,nomain,nonumberlist]{glossaries}
\makenoidxglossaries
%\usepackage{titlesec,titletoc} %turn on if tagging off and below
%%%
%\titleformat{\part}{\normalfont\large\bfseries\filcenter}{\thepart}{}{}[]
%\titlespacing*\part{0pt}{0.95\baselineskip}{0.50\baselineskip}
%\titleformat{\section}[runin]{\normalfont\normalsize\bfseries}{\thesection.}{0.5em}{}[]
%\titlespacing*\section{0pt}{0.50\baselineskip}{0.50\baselineskip}
%\titleformat{\subsection}[runin]{\normalfont\normalsize\bfseries}{\thesection\alph{subsection}.}{0.25em}{}[]
%\titlespacing*\subsection{0pt}{0.40\baselineskip}{0.50\baselineskip}
%\titleformat{\subsubsection}[runin]{\normalfont\normalsize\bfseries}{\roman{subsubsection}.}{0.25em}{}[]
%\titlespacing*\subsubsection{9pt}{0.25\baselineskip}{0.50\baselineskip}
%\titleformat{\paragraph}[runin]{\normalfont\normalsize\bfseries}{\theparagraph}{-1em}{}[]
%\titlespacing*\paragraph{0pt}{0.40\baselineskip}{0.50\baselineskip}
%\titleformat{\subparagraph}[runin]{\normalfont\normalsize\itshape}{\thesubparagraph}{0.25em}{}[]
%\titlespacing*\subparagraph{15pt}{0.25\baselineskip}{0.50\baselineskip}
%%%%%


%%%%% editing
\newcommand{\edit}[1]{\textcolor{blue}{#1}} %shortcut for changing font color on revised text
\newcommand{\fn}[1]{\footnote{#1}} %shortcut for footnote tag
\newcommand*\sq{\mathbin{\vcenter{\hbox{\rule{.3ex}{.3ex}}}}} %makes a small square as a separator $\sq$
%\newcommand{\sk}[1]{\sout{#1}} %shortcut for default strikethrough - do not sk through citep
\newcommand\sk{\bgroup\markoverwith{\textcolor{red}{\rule[0.5ex]{1pt}{1pt}}}\ULon} %strikethrough with red line; not in \section{}
%\st{} does strikethrough using soul package but does not like acronyms
\newcommand{\blucell}{\cellcolor{aliceblue}} %use to shade in table cell
\newcommand{\grycekk}{\cellcolor{lightgray}} %use to shade in table cell
\newcommand{\whicell}{\cellcolor{antiquewhite}} %use to shade in table cell
%%%%%


%%%%% colors
%http://latexcolor.com/
%https://en.wikibooks.org/wiki/LaTeX/Colors#:~:text=black%2C%20blue%2C%20brown%2C%20cyan,be%20available%20on%20all%20systems.
\definecolor{pridegold}{rgb}{241,179,0}
\definecolor{aliceblue}{rgb}{0.94, 0.97, 1.0}
\definecolor{antiquewhite}{rgb}{0.98, 0.92, 0.84}
\definecolor{lightmauve}{rgb}{0.86, 0.82, 1.0}
\definecolor{brilliantlavender}{rgb}{0.96, 0.73, 1.0}
\definecolor{brandeisblue}{rgb}{0.0, 0.44, 1.0}
\definecolor{darkmidnightblue}{rgb}{0.0, 0.2, 0.4}

\newcommand{\x}{\cellcolor{pridegold}} %use to shade in table cell
\newcommand{\y}{\cellcolor{lightgray}} %use to shade in table cell
\newcommand{\z}{\cellcolor{antiquewhite}} %use to shade in table cell
%%%%%


%%%%% acronyms
\newcommand{\acf}{\acrfull} %full acronym
\newcommand{\acl}{\acrlong} %long acronym
\newcommand{\acs}{\acrshort} %short acronym

\newcommand{\acfp}{\acrfullpl} %full acronym plural
\newcommand{\aclp}{\acrlongpl} %long acronym plural
\newcommand{\acsp}{\acrshortpl} %short acronym plural
%%%%%


%%%%% todonotes
\newcommandx{\cmt}[2][1=]{\todo[author=\textbf{STRUCTURE},tickmarkheight=0.15cm,linecolor=red,backgroundcolor=red!25,bordercolor=black,#1]{#2}}
\newcommandx{\con}[2][1=]{\todo[author=\textbf{CONTENT},tickmarkheight=0.15cm,linecolor=brilliantlavender,backgroundcolor=brilliantlavender,bordercolor=black,#1]{#2}}
\newcommandx{\rab}[2][1=]{\todo[noline,author=\textbf{RAB},backgroundcolor=Plum!25,bordercolor=black,#1]{#2}}


%\newcommandx{\jon}[2][1=]{\todo[noline,author=\textbf{ATTN: Johnson},backgroundcolor=blue!25,bordercolor=black,#1]{#2}}
%\newcommandx{\han}[2][1=]{\todo[noline,author=\textbf{ATTN: Haney},backgroundcolor=OliveGreen!25,bordercolor=black,#1]{#2}}
%\newcommandx{\rab}[2][1=]{\todo[author=\textbf{RAB},tickmarkheight=0.15cm,linecolor=Plum,backgroundcolor=Plum!25,bordercolor=black,#1]{#2}}
%\newcommandx{\han}[2][1=]{\todo[author=\textbf{ATTN: Haney},tickmarkheight=0.15cm,linecolor=OliveGreen,backgroundcolor=OliveGreen!25,bordercolor=OliveGreen,#1]{#2}}
%\newcommandx{\jon}[2][1=]{\todo[author=\textbf{ATTN: Johnson},tickmarkheight=0.15cm,linecolor=blue,backgroundcolor=blue!25,bordercolor=blue,#1]{#2}}


% highlighting 
\DeclareRobustCommand{\hlc}[1]{{\sethlcolor{LimeGreen}\hl{#1}}}
\makeatletter
    \if@todonotes@disabled
    \newcommand{\hlh}[2]{#1}
    \else
    \newcommand{\hlh}[2]{\han{#2}\hlc{#1}}
    \fi
    \makeatother

\DeclareRobustCommand{\hld}[1]{{\sethlcolor{CornflowerBlue}\hl{#1}}}
\makeatletter
    \if@todonotes@disabled
    \newcommand{\hlj}[2]{#1}
    \else
    \newcommand{\hlj}[2]{\jon{#2}\hld{#1}}
    \fi
    \makeatother

\DeclareRobustCommand{\hlf}[1]{{\sethlcolor{lightmauve}\hl{#1}}}
\makeatletter
    \if@todonotes@disabled
    \newcommand{\hlb}[2]{#1}
    \else
    \newcommand{\hlb}[2]{\rab{#2}\hlf{#1}}
    \fi
    \makeatother
%%%%%


%%%%% table alignments
\newcolumntype{L}[1]{>{\raggedright\let\newline\\\arraybackslash\hspace{0pt}}m{#1}} %uses \raggedright with m,p{} in table column
\newcolumntype{C}[1]{>{\centering\let\newline\\\arraybackslash\hspace{0pt}}m{#1}} %uses \raggedright with m,p{} in table column
\newcolumntype{R}[1]{>{\raggedleft\let\newline\\\arraybackslash\hspace{0pt}}m{#1}} %uses \raggedright with m,p{} in table column
%%%%%


%%%%% table contents
\makeatletter
\renewcommand\tableofcontents{%
    \@starttoc{toc}%
}
\makeatother

\makeatletter
\renewcommand\listoffigures{%
    \@starttoc{lof}%
}
\makeatother

\makeatletter
\renewcommand\listoftables{%
    \@starttoc{lot}%
}
\makeatother

\makeatletter
\newcommand*\ftp{\fontsize{16.5}{17.5}\selectfont}
\makeatother
%%%%%


%%%%% header and footer
\usepackage{fancyhdr}
\pagestyle{fancy}
\fancyhf{} %move page number to bottom right
%\renewcommand{\headrulewidth}{0pt} %set line thickness in header; uncomment as is to remove line
\lhead{\scriptsize NE585}
\chead{\scriptsize \textit{White Paper Project Proposal}}
\rhead{\scriptsize \today}
\rfoot{\thepage}
%%%%%


%%%%% acronyms
\input{$ACRONYM/acronyms}
%\input{../acronyms/acronyms}
%%%%%


%%%%% spacing
%\onehalfspacing %linespacing
%\setstretch{1.05} %linespacing
%\spacing{1.25} %equivalent to 1.5 line spacing in Word
%%%%%

%%%%% linenumbering
\linenumbers %toggle line numbers
\modulolinenumbers[5] %line numbering interval
%%%%%


\begin{document}


{\centering 
    \textbf{Title\\
    NE585 -- Nuclear Fuel Cycle Analysis Design Project\\
    }
    \vspace{0.10in}
    Name\\
    University of Idaho $\sq$ Idaho Falls Center for Higher Education\\
    Department of Nuclear Engineering and Industrial Management
\par
}

\noindent\makebox[\linewidth]{\rule{\textwidth}{0.5pt}} %horizontal line spanning margins


\section{Proposed Research}


\section{Motivation}


\section{Importance \& Relevance to Course Objectives}


\section{Background}
\subsection{Topic}

\subsection{Another topic}

\subsection{Related research}


\section{Methodology}
\subsection{One part of the methodology}

\subsection{Another part of the methodology}


\section{Major Outcomes \& Deliverables}


\begin{enumext}[series=outerlist,topsep=0pt,itemsep=-0.5ex,partopsep=1ex,parsep=1ex,label=\textbf{Workscope \arabic* --},align=left]
    \item\textbf{.} 
        \begin{enumext}[series=innerlist,topsep=3pt,itemsep=1.5ex,partopsep=1ex,parsep=-1ex,label=\textit{Task \Roman*.},labelsep=2.5ex,align=left,list-offset=-12ex]
            \item\textit{.}
            \item\textit{.}
        \end{enumext}
\end{enumext}

\begin{enumext}[series=milestone,topsep=0pt,itemsep=-0.5ex,partopsep=1ex,parsep=1ex,label=\textbf{\textit{Milestone \arabic*.}},align=left]
    \item\textit{.}
\end{enumext}


\begin{enumext}[resume=outerlist]
    \item\textbf{.}
        \begin{enumext}[resume=innerlist]
            \item\textit{.}
                \begin{enumext*}[topsep=0pt,itemsep=-0.5ex,partopsep=1ex,parsep=1ex,label=(\arabic*),columns=2]
                    \item .
                    \item .
                \end{enumext*}
            \item\textit{.}
                \begin{enumext*}[topsep=0pt,itemsep=-0.5ex,partopsep=1ex,parsep=1ex,label=(\arabic*),columns=2]
                    \item .
                    \item .
                    \item .
                    \item .
                \end{enumext*}
        \end{enumext}
\end{enumext}

\begin{enumext}[resume=milestone]
    \item\textit{.}
\end{enumext}


\begin{enumext}[resume=outerlist]
    \item\textbf{.}
        \begin{enumext}[resume=innerlist]
            \item\textit{.}
                \begin{enumext*}[topsep=0pt,itemsep=-0.5ex,partopsep=1ex,parsep=1ex,label=(\arabic*),columns=2]
                    \item .
                    \item .
                \end{enumext*}
            \item\textit{.}
                \begin{enumext*}[topsep=0pt,itemsep=-0.5ex,partopsep=1ex,parsep=1ex,label=(\arabic*),columns=2]
                    \item .
                    \item .
                    \item .
                    \item .
                \end{enumext*}
        \end{enumext}
\end{enumext}

\begin{enumext}[resume=milestone]
    \item\textit{.}
\end{enumext}


\begin{enumext}[resume=outerlist]
    \item\textbf{Crosscutting discussions \& Follow on activities.} 
        \begin{enumext}[resume=innerlist]
            \item\textit{Identify lessons learned.}
                \begin{enumext*}[topsep=0pt,itemsep=-0.5ex,partopsep=1ex,parsep=1ex,label=(\arabic*),columns=2]
                    \item .
                \end{enumext*}
            \item\textit{Develop future work.}
                \begin{enumext*}[topsep=0pt,itemsep=-0.5ex,partopsep=1ex,parsep=1ex,label=(\arabic*),columns=2]
                    \item .
                    \item . 
                    \item . 
                    \item .
                \end{enumext*}
            \item\textit{\acf{rv}.}
                \begin{enumext*}[topsep=0pt,itemsep=-0.5ex,partopsep=1ex,parsep=1ex,label=(\arabic*),columns=2]
                   \item Reproducibility -- .
                   \item Validation -- .
                \end{enumext*}
        \end{enumext}
\end{enumext}

\begin{enumext}[resume=milestone]
    \item\textit{Apply lessons learned for future research and engage colleagues in new collaborations.}
\end{enumext}


\section{Contribution to Advancing State of the Art}
This project will advance the state-of-the-art in nuclear fuel cycle analysis by --
\begin{enumext*}[topsep=0pt,itemsep=-0.5ex,partopsep=1ex,parsep=1ex,label=(\arabic*),columns=3]
    \item A; 
    \item B;
    \item C.
\end{enumext*}


\section{\textbf{Timeframe for Execution}}
\tagpdfsetup{table/header-rows=1}
\begin{longtable}{|L{0.17\linewidth}|c|c|c|c|c|}
        \caption{Project Timeline}
        \label{tab-timeline}
        \\
        \hline
        \multicolumn{1}{|c|}{\textbf{TASKS}}& 
        \textbf{AUG}& 
        \textbf{SEP}& 
        \textbf{OCT}& 
        \textbf{NOV}& 
        \textbf{DEC}
        \\
        \hline
        \endfirsthead
        I. 
        &\x
        & 
        & 
        & 
        & 
        \\
        \hline
        II. 
        &\x
        & 
        & 
        & 
        & 
        \\
        \hline
        III. 
        &\x 
        & 
        & 
        & 
        & 
        \\
        \hline
        IV. 
        &\x
        & 
        & 
        & 
        & 
        \\
        \hline
        V. 
        &\x 
        & 
        & 
        & 
        & 
        \\
        \hline
        VI. 
        &\x 
        & 
        & 
        & 
        & 
        \\
        \hline
        VII.
        & 
        & 
        & 
        & 
        &\x 
        \\
        \hline
\end{longtable}


\newpage


%%%%% header and footer
\pagestyle{fancyplain}
\fancyhf{} %clear header and footer 
\renewcommand{\headrulewidth}{0pt} %set line thickness in header; uncomment as is to remove line
\rfoot{\fancyplain{}{\thepage}}
\nolinenumbers
%%%%%


\bibliographystyle{nsf}
\setlength{\bibhang}{0pt}
\bibliography{}


\newpage


\begin{appendices}


%%%%% formatting
    \renewcommand{\thesection}{\Roman{section}}
%%%%%


\newpage


    \section{Acronyms}
    \printnoidxglossary[title={}]


\end{appendices}


%wrap figure around text
%move to wherever in the text 
%{l} fixes it on the left margin
%{L} floats
%\begin{wrapfigure}{l}{0.x\textwidth}
%       \includegraphics[width=0.(x-2)\textwidth]{}
%    \caption{}
%    \label{fig-}
%\end{wrapfigure}


\end{document}
